\chapter{Discussion}

Overall, the research met its three goals of performance, modularity and portability. The GPU
backing made large simulations substantially faster, custom functions let researchers describe new
pulses and couplings without changing code, and the simulation still builds and runs on machines
without a GPU. Modern C++ features proved useful in getting there. Holding each backing in a
\texttt{std::variant} and dispatching with \texttt{std::visit} kept a single propagation path for both
the CPU and GPU, and meant that a missing implementation was caught by the compiler rather than at
runtime. At the same time, the new code kept to the conventions already in the codebase, such as
configuring components through separate \texttt{setup()} calls. While this departs from standard C++
practice, it is common in scientific computing and kept the new work consistent with how the rest of
the simulation is built.

There are several directions for further work. The single-precision results were promising on speed
but introduced too much error, so more experimentation with 32-bit values, such as keeping only the
FFTs in double precision, could recover much of the speedup without the loss in accuracy. AMD GPUs
are not yet supported, and adding and testing ROCm support would make GPU acceleration available to
more researchers. Adding another \texttt{compute\_backing} type for a library such as SYCL
would give lower-level control over GPU computation, and show that \texttt{FieldData} can support new
data storage.
Finally, the performance improvements were measured on benchmark configurations, and demonstrating them 
on an actual research question would show how much they shorten the time taken to get real results.
